% Vietnamese translation for OI Public Library
% Source-File: IOI中国国家候选队论文/2008/Day2/2.高逸涵《部分贪心思想在信息学竞赛中的应用》/部分贪心在信息学竞赛中的应用.ppt
\documentclass[11pt]{article}
\usepackage{oipl-vn}

\newcommand{\Oh}{\mathcal{O}}

\title{Bản trình chiếu: Ứng dụng tư tưởng tham lam một phần}
\author{Gao Yihan}
\date{2008}

\begin{document}
\maketitle

\origtitle{Slide companion for partial greedy thinking}
\sourcepath{IOI China candidate team papers / 2008 / Day 2 / Gao Yihan.ppt}

\section{Phạm vi bản dịch}

Tệp PowerPoint đi cùng bài viết đã được dịch từ tài liệu Word. Bản ghi chú này
giữ lại ý chính: dùng tham lam ở phần chắc chắn đúng để thu nhỏ bài toán, rồi
dùng thuật toán chắc chắn hơn cho phần còn lại.

\section{Vì sao cần tham lam một phần}

Tham lam thuần túy nhanh và dễ viết, nhưng thường có phản ví dụ nhỏ. Quy hoạch
động hoặc tìm kiếm đáng tin hơn, nhưng có thể không chịu nổi dữ liệu lớn. Tham
lam một phần chọn một vùng an toàn: khi quy mô còn lớn, thực hiện bước tham
lam có thể chứng minh; khi chỉ còn phần nhỏ hoặc phần biên, chuyển sang liệt
kê, DP hoặc tìm kiếm.

\section{Mẫu sử dụng}

Trong các ví dụ của bài viết, tham lam dùng để đưa nghiệm tới gần vùng tối ưu,
hoặc để giảm dữ liệu xuống dưới một ngưỡng. Sau đó phần còn lại được giải
chính xác. Cách này giữ được tốc độ gần như tham lam, nhưng tránh lỗi do các
trường hợp biên nhỏ.

\section{Kết luận}

Tham lam một phần không phải là dùng trực giác thay chứng minh. Điều cần chứng
minh là sau mỗi bước tham lam, nghiệm tối ưu vẫn nằm trong không gian còn lại,
hoặc phần bị bỏ đi không thể ảnh hưởng tới đáp án cuối cùng.

\end{document}
